% ============================================================
%  MODULE IX — Macro & Global Markets
% ============================================================

\chapter{Macro \& Global Markets}

\begin{notebox}
No asset exists in isolation. Stock prices, bond yields, exchange rates, and commodity prices
are connected through \textbf{macroeconomic forces}. Understanding the macro backdrop is essential
for context — even if you trade purely technically.
\end{notebox}

% ────────────────────────────────────────────────────────────
\section{The Economic Cycle}

The economy moves in cycles: \textbf{expansion} $\rightarrow$ \textbf{peak} $\rightarrow$
\textbf{contraction/recession} $\rightarrow$ \textbf{trough} $\rightarrow$ expansion again.

\subsection{Sector Performance Across the Cycle}

\begin{center}
\begin{tabular}{>{\bfseries}p{3cm} p{4cm} p{6cm}}
\toprule
Cycle Phase & Economy Characteristics & Outperforming Sectors \\
\midrule
Early Expansion & Recovery, low rates, rising confidence & Consumer Discretionary, Financials \\
Mid Expansion & Solid growth, moderate inflation & Technology, Industrials, Materials \\
Late Expansion & High inflation, rising rates & Energy, Commodities, Real Estate \\
Recession & Contraction, falling rates & Consumer Staples, Utilities, Health Care \\
\bottomrule
\end{tabular}
\end{center}

% ────────────────────────────────────────────────────────────
\section{Interest Rates and the Bond Market}

\subsection{Central Banks and the Policy Rate}
The central bank sets a \textbf{short-term policy interest rate} (e.g., the Federal Funds Rate).
This rate propagates through the entire financial system:

\begin{itemize}
  \item \textbf{Higher rates} → borrowing more expensive → lower business investment, lower
        consumer spending → slows inflation but also growth
  \item \textbf{Lower rates} → borrowing cheaper → more investment, more consumption →
        stimulates growth but risks inflation
\end{itemize}

\subsection{The Yield Curve}
A plot of \textbf{bond yields} (interest rates) across maturities (3m, 1y, 2y, 5y, 10y, 30y).

\begin{keyterm}{Yield Curve Inversion}
When short-term yields exceed long-term yields. Historically, an inverted yield curve
has preceded every US recession since the 1950s. The 2-year vs 10-year spread is most watched.
\end{keyterm}

\subsection{Bond Prices and Yields — Inverse Relationship}

\begin{formulabox}
\[
  \text{Bond Price} \propto \frac{1}{\text{Yield}}
\]
When yields rise, existing bond prices fall. When yields fall, bond prices rise.
This is because the fixed coupon payment becomes relatively less/more attractive
compared to new bonds.
\end{formulabox}

\subsection{How Rates Affect Stock Valuations}
Higher interest rates:
\begin{enumerate}[label=\arabic*.]
  \item Raise the discount rate in DCF models → lower present value → lower stock prices.
  \item Make risk-free returns (bonds, cash) more competitive vs. stocks.
  \item Increase corporate borrowing costs → compress earnings.
\end{enumerate}

Growth stocks (with cash flows far in the future) are especially rate-sensitive.

% ────────────────────────────────────────────────────────────
\section{Inflation}

\begin{keyterm}{Inflation}
The rate at which the general level of prices for goods and services rises,
eroding purchasing power. Measured by CPI (Consumer Price Index) and PCE (Personal
Consumption Expenditures).
\end{keyterm}

\subsection{Inflation's Effect on Markets}
\begin{itemize}
  \item \textbf{Moderate inflation} (2--3\%): Healthy; supports nominal earnings growth.
  \item \textbf{High inflation}: Central banks raise rates → negative for bonds and growth stocks;
        commodities and energy often outperform.
  \item \textbf{Deflation}: Very negative; reduces corporate revenues; Japan's ``Lost Decades''.
\end{itemize}

% ────────────────────────────────────────────────────────────
\section{Key Macro Economic Indicators}

\begin{center}
\begin{tabular}{>{\bfseries}p{3.5cm} p{3cm} p{6.5cm}}
\toprule
Indicator & Release Frequency & What Traders Watch \\
\midrule
Non-Farm Payrolls (NFP) & Monthly (US) & Employment health; market-moving \\
CPI / Core CPI & Monthly & Inflation trajectory; Fed direction \\
GDP Growth & Quarterly & Economic expansion/contraction \\
PMI (Manufacturing/Services) & Monthly & Real-time business activity \\
Retail Sales & Monthly & Consumer spending health \\
Jobless Claims & Weekly & Leading employment indicator \\
Fed Funds Rate Decision & 8x/year & Defines cost of money \\
FOMC Minutes & 8x/year & Fed deliberation detail \\
\bottomrule
\end{tabular}
\end{center}

\begin{warningbox}
Major macro releases (\textbf{NFP}, \textbf{CPI}, \textbf{FOMC}) cause violent, fast price moves
with wide spreads. Experienced traders often close or hedge positions before these releases.
Never hold leveraged positions through major macro events unless it is part of a deliberate,
sized strategy.
\end{warningbox}

% ────────────────────────────────────────────────────────────
\section{Currency Markets (FX)}

Exchange rates affect:
\begin{itemize}
  \item \textbf{Multinational earnings} — a strong USD reduces the dollar value of foreign revenues.
  \item \textbf{Commodity prices} — most commodities are priced in USD; strong dollar → lower prices.
  \item \textbf{Emerging markets} — EM countries with USD-denominated debt suffer when USD strengthens.
\end{itemize}

\subsection{Key Currency Relationships}
\begin{itemize}
  \item \textbf{USD/JPY} — risk barometer; yen strengthens in risk-off environments
  \item \textbf{AUD/USD, CAD/USD} — commodity currencies; follow commodity prices
  \item \textbf{EUR/USD} — most liquid FX pair; follows ECB vs Fed policy divergence
  \item \textbf{DXY} — US Dollar Index; inverse relationship with most risk assets
\end{itemize}

% ────────────────────────────────────────────────────────────
\section{Commodities}

\begin{center}
\begin{tabular}{>{\bfseries}p{3cm} p{4cm} p{6cm}}
\toprule
Commodity & Category & Key Drivers \\
\midrule
Crude Oil (WTI/Brent) & Energy & OPEC policy, global demand, USD \\
Gold & Precious Metal & Real interest rates, USD, safe haven demand \\
Silver & Precious/Industrial & Gold correlation + industrial demand \\
Copper & Industrial Metal & China demand; economic growth proxy \\
Wheat/Corn/Soy & Agriculturals & Weather, geopolitics, USD \\
\bottomrule
\end{tabular}
\end{center}

\textbf{Gold and real interest rates} have a strong inverse relationship:
\[
  \text{Real Rate} = \text{Nominal Rate} - \text{Expected Inflation}
\]
When real rates fall (e.g., negative real rates), gold becomes more attractive as a store
of value vs. bonds.

% ────────────────────────────────────────────────────────────
\section{Inter-Market Analysis}

The four major asset classes — equities, bonds, commodities, and currencies — are
linked. Classic relationships (not always reliable):

\begin{itemize}
  \item Bonds $\uparrow$ (yields $\downarrow$) $\rightarrow$ stocks $\uparrow$ (historically)
  \item Dollar $\uparrow$ $\rightarrow$ gold $\downarrow$, commodities $\downarrow$, EM $\downarrow$
  \item Commodities $\uparrow$ $\rightarrow$ inflation expectations $\uparrow$ $\rightarrow$ bonds $\downarrow$
  \item Copper $\uparrow$ $\rightarrow$ global growth optimism $\rightarrow$ cyclical stocks $\uparrow$
\end{itemize}

\begin{notebox}
These correlations \textbf{shift through regimes}. The bond-equity correlation was negative for
20 years (2000--2020) and turned positive in the 2022 inflation shock. Always verify
whether historical relationships still hold in the current regime.
\end{notebox}
